This appendix does what the four-page body has no room for: it places the
episode of Section \ref{sec:t9} against the work it resembles, and places the
mathematical substrate of Section \ref{sec:substrate} against the literatures
it draws on. Neither placement changes a claim in the body.

\subsection{Machines and people in mathematics}

The closest precedent for the shape of our loop is the use of computation to
guide mathematical intuition toward statements people then prove
\citep{davies2021advancing}, where a learned model's saliency over a space of
invariants pointed at the relationships that became theorems in knot theory
and representation theory. The direction of travel there is machine to human:
the computation supplies a candidate, the person supplies the proof. Our
eleventh pivot runs the other way for one leg of the cycle, since the object
the proof supplied was not a candidate statement but a region, and the
region is what made the second search productive.

A second line lets search refute conjectures directly. Reinforcement learning
over constructions has produced explicit counterexamples in extremal
combinatorics and spectral graph theory \citep{wagner2021constructions}, and
program search guided by a language model has produced both counterexample-like
constructions and improved bounds \citep{romera2024mathematical}. Our search is
far less sophisticated than either, and that is the point of the episode: the
first search was not underpowered in any way its own diagnostics could see. It
was thorough over a space that excluded the answer, and it reported a minimum
of $1.005$ with two independent strategies agreeing. What told us the space was
wrong was a proof that stopped at $\hstar - 2/c$ and said so.

A third line automates proof itself, from olympiad geometry
\citep{trinh2024solving} to olympiad-level formal reasoning with reinforcement
learning \citep{hubert2025olympiad}, on the infrastructure of large formal
libraries \citep{mathlib2020}. We do none of this. Our proofs are written and
checked by people and are not formalized, which is a limitation we state rather
than defend. The older lineage of machine-assisted proof, from the four-colour
theorem \citep{appel1977every} to the formal Kepler proof
\citep{hales2017formal}, and the experimental-mathematics tradition of treating
computation as evidence to be believed provisionally
\citep{borwein2008mathematics}, is the tradition our verification harness sits
in: it is bookkeeping about how much a numerical verdict is worth, not proof.

What we claim to add is narrow. It is a completely recorded instance of the
failure direction, in which a numerical verdict was written down as a result,
was wrong, and was corrected by a proof that could not reach it. The record
includes the search that produced the false verdict, the reason it missed
(a weight grid floored at $0.05$ against a violating probe of
$3 \times 10^{-4}$, and a symmetric-pair prior the violating design does not
satisfy), the proof's exact stopping point, and the frozen witness.

\subsection{The substrate's literatures}

\textbf{Performative prediction.} The fixed-point framing and the gap between
the stable point and the optimum are due to \citet{perdomo2020performative},
surveyed in \citet{hardt2023performative}. Algorithms that close the gap
\citep{izzo2021learn,miller2021outside,mendler2020stochastic,drusvyatskiy2022stochastic}
and their stateful variants \citep{brown2022performative,li2022state} all
consume an estimate of the distribution's response, obtained by perturbing
deployments, and treat those perturbations as free.
\citet{jagadeesan2022regret} bound the regret of \emph{finding} the
performative optimum; the exchange rate prices \emph{knowing your own
feedback} instead, giving lower bounds and an invariance identity where they
give algorithmic upper bounds. \citet{bracale2025reverse} design deployments
to learn the distribution map and bound the regret of doing so at an order
level; T2 prices the same activity exactly, in the objective's own curvature.
\citet{lin2024plugin} show that misspecified structural models still help,
which T7 sharpens into a budget-explicit rule.

\textbf{Learning while earning.} That myopic optimization can stall learning is
classical in operations research and economics
\citep{rothschild1974two,easley1988controlling,aghion1991optimal,harrison2012bayesian,broder2012dynamic,keskin2014dynamic},
surveyed by \citet{denboer2015dynamic}, with the informational tension stated
in its sharpest form by \citet{grossman1980impossibility}. T1 is the
performative version of that stall: it comes from the optimizer converging,
not from weak feedback, and the cap is explicit. \citet{lai1979adaptive}'s
adaptive designs are points on the frontier T2 describes; what does not appear
in that literature is the frontier itself, its floor, or a demand curve the
decision-maker's own actions generate.

\textbf{Dual control and experiment design.} Pricing an experiment inside the
objective it serves is dual control \citep{feldbaum1960dual} and least-costly
identification \citep{bombois2006least}. Here the plant is a retraining
learner, the budget is the optimizer's own objective, and the constraint is
loop contraction. The lower bounds are van Trees arguments
\citep{gill1995applications}. On the design side, task-optimal exploration in
linear systems \citep{wagenmaker2021task,dean2020sample} and the finding that
naive exploration suffices for the linear quadratic regulator
\citep{simchowitz2020naive} sit against C5.1, which says isotropic exploration
overpays by an explicit curvature-dispersion factor once the design is priced
in the experimenter's own objective. The classical canon
\citep{kiefer1960equivalence,pukelsheim2006optimal} supplies the A-optimality
and Chebyshev-system machinery but never prices the design that way.

\textbf{The market.} The structural single-book market, its calibration, and
the corporate-bond instantiation behind the real-data row of Table
\ref{tab:results} are those of REFLEX \citep{reflex2026}. The
quote-summons-toxicity channel is the dealer's adverse-selection problem
\citep{gueant2013inventory}, and the fact that a dealer's own quotes feed the
flow it then reads is the price-reading effect studied by
\citet{barzykin2026quoting}.
